\section{Task 2B: Dynamic I-RRIP with Set Dueling}
\label{sec:task2b}
\subsection{Set selection and decision statistic}
The supplied \texttt{dyn\_irrip.l2c\_repl} sets \texttt{maxRRPV=3}, uses \texttt{RESET\_INTERVAL=4096}, and defines the threshold as the compile-time constant \texttt{THRESHOLD 0.25}. For set index $s$, sets with $s\bmod32=0$ are RRIP leader sets, sets with $s\bmod32=1$ are I-RRIP leader sets, and all other sets use the global follower decision. The two leader classes together account for 2 of every 32 sets (6.25\% when the number of sets is divisible by 32), with each leader class representing 3.125\% of sets.

The policy updates shared counters on accesses to either leader class and computes
\[
R=\frac{A_I+M_I}{A_{\mathrm{total}}+M_{\mathrm{total}}},
\]
where the \texttt{instr\_accesses}, \texttt{instr\_misses}, \texttt{total\_accesses}, and \texttt{total\_misses} counters implement these quantities. The ratio is recalculated after each leader-set callback and compared with 0.25; follower sets use I-RRIP when $R\geq0.25$ and RRIP otherwise. Once total accesses reach 4,096, the counters are reset after the ratio update. The source uses a single shared statistic rather than keeping independent utility counters for the two leader policies, so its decision is instruction-activity based rather than a direct comparison of measured miss cost between the two leader groups.

\subsection{Threshold sensitivity: a source-level issue}
The runner attempted thresholds 0.10, 0.20, 0.25, 0.30, and 0.40. The measured IPC values and cache metrics repeat exactly across those five rows. This is explained by the implementation: \texttt{dyn\_irrip.l2c\_repl} hardcodes \texttt{THRESHOLD 0.25} and never reads the \texttt{CHAMPSIM\_THRESHOLD} environment variable that the runner changes. Therefore these runs all use the same effective threshold and do \emph{not} constitute a valid threshold sensitivity experiment. The repeated table below preserves the observed results for reproducibility; it should not be used to claim that the policy is insensitive to threshold or that 0.25 is optimal.

\begin{table}[H]\centering\small
\caption{Observed Dynamic I-RRIP IPC by requested threshold. The invariance is an artifact of the compile-time constant, not evidence from a working sweep.}
\begin{tabular}{lrrrr}
\toprule Requested threshold & Tr1 IPC & Tr2 IPC & Tr3 IPC & Mean speedup vs RRIP \\
\midrule
0.10 & 0.232281 & 0.256285 & 0.381742 & 1.006004 \\
0.20 & 0.232281 & 0.256285 & 0.381742 & 1.006004 \\
0.25 & 0.232281 & 0.256285 & 0.381742 & 1.006004 \\
0.30 & 0.232281 & 0.256285 & 0.381742 & 1.006004 \\
0.40 & 0.232281 & 0.256285 & 0.381742 & 1.006004 \\
\bottomrule\end{tabular}\end{table}
\plotpair{plots/task2b/sensitivity_ipc.png}{IPC by requested threshold}{plots/task2b/sensitivity_l2i_load_mpki.png}{L2 instruction MPKI by requested threshold}{Requested-threshold sweep; identical bars reflect the hardcoded threshold described above}{fig:t2b-sensitivity}

\subsection{Measured policy comparison at effective threshold 0.25}
\begin{table}[H]\centering\scriptsize
\caption{Task 2B results. Speedup is relative to RRIP on the same trace.}
\resizebox{\linewidth}{!}{\begin{tabular}{llrrrrr}
\toprule Trace & Policy & IPC & Speedup & L2C Load MPKI & L2 Instruction MPKI & L2 Data MPKI \\
\midrule
Tr1.gz & RRIP baseline & 0.232231 & 1.000000 & 2.2021 & 10.5618 & 7.0769 \\
Tr1.gz & I-RRIP & 0.232164 & 0.999711 & 2.4089 & 9.8916 & 7.7172 \\
Tr1.gz & Dynamic I-RRIP (compiled threshold 0.25) & 0.232281 & 1.000215 & 2.3941 & 9.9094 & 7.6956 \\
Tr2.gz & RRIP baseline & 0.251806 & 1.000000 & 2.1711 & 17.7813 & 7.9645 \\
Tr2.gz & I-RRIP & 0.256625 & 1.019138 & 2.3754 & 13.0364 & 9.1737 \\
Tr2.gz & Dynamic I-RRIP (compiled threshold 0.25) & 0.256285 & 1.017788 & 2.3631 & 13.3453 & 9.1221 \\
Tr3.gz & RRIP baseline & 0.381739 & 1.000000 & 0.7090 & 1.1970 & 2.8424 \\
Tr3.gz & I-RRIP & 0.381981 & 1.000634 & 0.7264 & 1.0221 & 2.8761 \\
Tr3.gz & Dynamic I-RRIP (compiled threshold 0.25) & 0.381742 & 1.000008 & 0.7235 & 1.1464 & 2.8764 \\
\bottomrule\end{tabular}}
\end{table}
\plotpair{plots/task2b/ipc.png}{IPC}{plots/task2b/speedup.png}{Speedup}{Task 2B policy-level IPC and speedup}{fig:t2b-ipc-speedup}
\plotpair{plots/task2b/l2c_load_mpki.png}{L2C load MPKI}{plots/task2b/l2i_load_mpki.png}{L2 instruction-load MPKI}{Task 2B aggregate and instruction MPKI}{fig:t2b-l2c-l2i}
\plotpair{plots/task2b/l2d_load_mpki.png}{L2 data-load MPKI}{plots/task2b/ipc_change_vs_rrip_by_threshold_pct.png}{IPC change vs. RRIP (\%)}{Task 2B data-MPKI trade-off and relative IPC change}{fig:t2b-l2d-change}
\plotone{plots/task2b/l2i_mpki_reduction_vs_rrip_by_threshold_pct.png}{Instruction-load MPKI reduction relative to RRIP for the dynamic policy at the reported threshold.}{fig:t2b-l2i-reduction}

\subsection{Interpretation and required correction}
At the measured effective threshold of 0.25, dynamic I-RRIP changes IPC by +0.0215\% on Tr1, +1.7788\% on Tr2, and approximately +0.0008\% on Tr3 relative to RRIP. L2 instruction MPKI declines on all three traces, while aggregate L2C and data MPKI rise. Tr2 is the clearest performance win in the supplied data; Tr1 and Tr3 are essentially neutral in IPC.

To complete the assignment's sensitivity requirement, replace the compile-time macro with a runtime parameter read from \texttt{CHAMPSIM\_THRESHOLD} (with a safe default and range check), rebuild for every threshold or ensure the runtime argument reaches the policy, and rerun all three traces at the five values. Keep leader-set placement and all other parameters fixed. Until those measurements exist, no threshold selection rationale can be inferred from this sweep.
